% Modèle matOS : notes de cours. Compiler avec pdflatex (TeXstudio : F5).
\documentclass[11pt,a4paper]{article}
\usepackage[T1]{fontenc}
\usepackage[utf8]{inputenc}
\usepackage[french]{babel}
\usepackage{lmodern,microtype}
\usepackage[margin=2.5cm]{geometry}
\usepackage{amsmath,amssymb,amsthm,mathtools}
\usepackage{enumitem}
\usepackage{xcolor}
\usepackage[colorlinks,linkcolor=blue!50!black]{hyperref}

\theoremstyle{plain}
\newtheorem{theoreme}{Théorème}[section]
\newtheorem{proposition}[theoreme]{Proposition}
\newtheorem{lemme}[theoreme]{Lemme}
\newtheorem{corollaire}[theoreme]{Corollaire}
\theoremstyle{definition}
\newtheorem{definition}[theoreme]{Définition}
\newtheorem{exemple}[theoreme]{Exemple}
\theoremstyle{remark}
\newtheorem{remarque}[theoreme]{Remarque}

\newcommand{\N}{\mathbb{N}}
\newcommand{\Z}{\mathbb{Z}}
\newcommand{\Q}{\mathbb{Q}}
\newcommand{\R}{\mathbb{R}}
\newcommand{\C}{\mathbb{C}}

\title{Titre du cours}
\author{Nom de l'enseignant \\ \small Mention Mathématiques, niveau L1}
\date{Année universitaire 2026--2027}

\begin{document}
\maketitle
\tableofcontents

\section{Suites réelles}

\begin{definition}
Une suite réelle $(u_n)_{n \in \N}$ \emph{converge} vers $\ell \in \R$ si
\[
  \forall \varepsilon > 0,\ \exists N \in \N,\ \forall n \geq N,\quad |u_n - \ell| \leq \varepsilon .
\]
\end{definition}

\begin{theoreme}[Unicité de la limite]
Si une suite réelle converge, sa limite est unique.
\end{theoreme}

\begin{proof}
Supposons que $(u_n)$ converge vers $\ell$ et vers $\ell'$. Pour tout $\varepsilon > 0$ et
tout $n$ assez grand, $|\ell - \ell'| \leq |\ell - u_n| + |u_n - \ell'| \leq 2\varepsilon$.
Donc $\ell = \ell'$.
\end{proof}

\begin{exemple}
La suite $u_n = \dfrac{1}{n+1}$ converge vers $0$.
\end{exemple}

\end{document}
